\chapter{2014年山东卷（文）}

\section{单选题}

\begin{problem}
已知 \(a\)，\(b \in \R\)，\(\i\) 是虚数单位，若 \(a + \i = 2 - b\i\)，则 \((a + b\i)^2 =\)（\quad）

\choices
    {\(3 - 4\i\)}
    {\(3 + 4\i\)}
    {\(4 - 3\i\)}
    {\(4 + 3\i\)}
\end{problem}

\begin{answer}
A
\end{answer}

\begin{solution}
由复数相等的充要条件，比较等式两边的实部和虚部，得 \(a=2\)，\(-b=1\)，所以 \(b=-1\). 因而
\[
(a+b\i)^2=(2-\i)^2=4-4\i+\i^2=3-4\i
\]
故选 A.
\end{solution}

\begin{problem}
设集合 \(A = \{x \mid x^2 - 2x < 0\}\)，\(B = \{x \mid 1 \leqslant x \leqslant 4\}\)，则 \(A \cap B =\)（\quad）

\choices
    {\((0,2]\)}
    {\( (1,2) \)}
    {\([1, 2)\)}
    {\((1,4)\)}
\end{problem}

\begin{answer}
C
\end{answer}

\begin{solution}
由 \(x^2-2x=x(x-2)<0\)，得 \(0<x<2\)，所以 \(A=(0,2)\). 又 \(B=[1,4]\)，从而
\[
A\cap B=[1,2)
\]
故选 C.
\end{solution}

\begin{problem}
函数 \( f(x) = \frac{1}{\sqrt{\log_2 x - 1}} \) 的定义域为（\quad）

\choices
    {\((0,2)\)}
    {\((0,2]\)}
    {\((2, + \infty)\)}
    {\([2, +\infty)\)}
\end{problem}

\begin{answer}
C
\end{answer}

\begin{solution}
根式在分母中，要求根式内的式子严格大于零，因此
\[
\log_2x-1>0
\]
由于底数 \(2>1\)，上式等价于 \(\log_2x>1\)，即 \(x>2\). 所以函数的定义域为 \((2,+\infty)\)，故选 C.
\end{solution}

\begin{problem}
用反证法证明命题“设 \(a\)，\(b\) 为实数，则方程 \(x^3 + ax + b = 0\) 至少有一个实根”时，要做的假设是（\quad）

\choices
    {方程 \( x^{3} + ax + b = 0 \) 没有实根}
    {方程 \(x^{3} + ax + b = 0\) 至多有一个实根}
    {方程 \(x^{3} + ax + b = 0\) 至多有两个实根}
    {方程 \(x^{3} + ax + b = 0\) 恰好有两个实根}
\end{problem}

\begin{answer}
A
\end{answer}

\begin{solution}
原命题为“方程至少有一个实根”. 它的否定是“方程没有实根”，所以用反证法时应假设方程
\[
x^3+ax+b=0
\]
没有实根.
故选 A.
\end{solution}

\begin{problem}
已知实数 \(x\)，\(y\) 满足 \(a^x < a^y\) （\(0 < a < 1\)），则下列关系式恒成立的是（\quad）

\choices
    {\(x^3 > y^3\)}
    {\(\sin x > \sin y\)}
    {\(\ln (x^2 + 1) > \ln (y^2 + 1)\)}
    {\(\frac{1}{x^2 + 1} >\frac{1}{y^2 + 1}\)}
\end{problem}

\begin{answer}
A
\end{answer}

\begin{solution}
当 \(0<a<1\) 时，函数 \(a^x\) 在 \(\R\) 上单调递减. 由 \(a^x<a^y\) 可得 \(x>y\). 由于函数 \(x^3\) 在 \(\R\) 上单调递增，所以
\[
x^3>y^3
\]
故选 A.
\end{solution}

\begin{problem}
已知函数 \( y = \log_a (x + c) \) （\(a\)，\(c\) 为常数，其中 \(a > 0\)，\(a \neq 1\)）的图象如图，则下列结论成立的是（\quad）

\bitmapfigure[max width=0.30\linewidth]{img/2014/shandong_liberal/q6_fig1.png}

\choices
    {\(a > 1, c > 1\)}
    {\(a > 1\)，\(0 < c < 1\)}
    {\(0 < a < 1\)，\(c > 1\)}
    {\(0 < a < 1\)，\(0 < c < 1\)}
\end{problem}

\begin{answer}
D
\end{answer}

\begin{solution}
由图象可知函数单调递减，因此 \(0<a<1\). 图象与 \(x\) 轴的交点横坐标在 \(0\) 与 \(1\) 之间，而当 \(y=0\) 时，
\(\log_a(x+c)=0\)，\(\Longrightarrow\)，\(x=1-c\).
所以 \(0<1-c<1\)，即 \(0<c<1\). 故选 D.
\end{solution}

\begin{problem}
已知向量 \(\bs{a} = (1, \sqrt{3})\)，\(\bs{b} = (3, m)\). 若向量 \(\bs{a}\)，\(\bs{b}\) 的夹角为 \(\frac{\pi}{6}\)，则实数 \(m =\)（\quad）

\choices
    {\(2\sqrt{3}\)}
    {\(\sqrt{3}\)}
    {\(0\)}
    {\(-\sqrt{3}\)}
\end{problem}

\begin{answer}
B
\end{answer}

\begin{solution}
由向量夹角公式，得
\[
\cos\frac\pi6=\frac{\bs a\cdot\bs b}{|\bs a||\bs b|}
 =\frac{3+\sqrt3m}{2\sqrt{9+m^2}}
\]
于是
\[
3+\sqrt3m=\sqrt3\sqrt{9+m^2}
\]
两边平方并化简，得
\(9+6\sqrt3m+3m^2=27+3m^2\)，\(m=\sqrt3\).
代入原等式可知等式成立，故选 B.
\end{solution}

\begin{problem}
为了研究某种药品的疗效，选取若干名志愿者进行临床试验. 所有志愿者的舒张压数据（单位：\(\mathrm{kPa}\)）的分组区间为 \([12,13)\)，\([13,14)\)，\([14,15)\)，\([15,16)\)，\([16,17]\). 将其按从左到右的顺序分别编号为第一组，第二组，…，第五组. 如图所示是根据试验数据制成的频率分布直方图. 已知第一组与第二组共有 \(20\text{ 人}\)，第三组中没有疗效的有 \(6\text{ 人}\)，则第三组中有疗效的人数为（\quad）

\texfigure[max width=0.34\linewidth]{img_repaint/2014/shandong_liberal/q8_fig1.tex}

\choices
    {\(6\)}
    {\(8\)}
    {\(12\)}
    {\(18\)}
\end{problem}

\begin{answer}
C
\end{answer}

\begin{solution}
频率分布直方图纵轴表示频率与组距的比值，而五组的组距均为 \(1\)，所以第一组、第二组和第三组的频率分别为 \(0.24\)、\(0.16\) 和 \(0.36\). 设志愿者总人数为 \(N\)，则
\[
(0.24+0.16)N=20
\]
解得 \(N=50\). 因此第三组共有 \(0.36\times50=18\) 人，其中有疗效的人数为
\[
18-6=12
\]
故选 C.
\end{solution}

\begin{problem}
对于函数 \( f(x) \)，若存在常数 \( a \neq 0 \)，使得 \( x \) 取定义域内的每一个值，都有 \( f(x) = f(2a - x) \)，则称 \( f(x) \) 为准偶函数，下列函数中是准偶函数的是（\quad）

\choices
    {\(f(x) = \sqrt{x}\)}
    {\(f(x) = x^{3}\)}
    {\(f(x) = \tan x\)}
    {\(f(x) = \cos (x + 1)\)}
\end{problem}

\begin{answer}
D
\end{answer}

\begin{solution}
取 \(a=-1\)，则 \(a\ne0\)，且对定义域内任意的 \(x\)，有
\[
f(2a-x)=\cos(-2-x+1)=\cos(-x-1)=\cos(x+1)=f(x)
\]
所以 \(f(x)=\cos(x+1)\) 是准偶函数，故选 D.
\end{solution}

\begin{problem}
已知 \(x\)，\(y\) 满足约束条件 \(\begin{cases}
x - y - 1 \leqslant 0,\\
2x - y - 3 \geqslant 0,
\end{cases}\) 当目标函数 \(z = ax + by\)（\(a > 0\)，\(b > 0\)）在该约束条件下取到最小值 \(2\sqrt{5}\) 时，\(a^2 + b^2\) 的最小值为（\quad）

\choices
    {\(5\)}
    {\(4\)}
    {\(\sqrt{5}\)}
    {\(2\)}
\end{problem}

\begin{answer}
B
\end{answer}

\begin{solution}
约束条件等价于
\(y\geqslant x-1\)，\(y\leqslant2x-3\).
要使可行域非空，必须有 \(x-1\leqslant2x-3\)，所以 \(x\geqslant2\). 对固定的 \(x\)，因为 \(b>0\)，目标函数在最小的可行 \(y\)，即 \(y=x-1\) 处取最小值. 因而
\[
z=ax+b(x-1)=(a+b)x-b
\]
又因为 \(a+b>0\) 且 \(x\geqslant2\)，所以目标函数的最小值在 \((x,y)=(2,1)\) 处取得，且
\[
2a+b=2\sqrt5
\]
由柯西不等式，
\[
(2a+b)^2\leqslant(2^2+1^2)(a^2+b^2)=5(a^2+b^2)
\]
所以 \(a^2+b^2\geqslant4\). 当 \((a,b)=t(2,1)\)，且 \(5t=2\sqrt5\)，即
\(a=\frac{4\sqrt5}{5}\)，\(b=\frac{2\sqrt5}{5}\)
时等号成立. 因此 \(a^2+b^2\) 的最小值为 \(4\)，故选 B.
\end{solution}

\section{填空题}

\begin{problem}
执行如图所示的程序框图，若输入的 \(x\) 的值为 1，则输出的 \(n\) 的值为\fillinblank{}.
\parbox{\dimexpr\linewidth-\problembodyindent\relax}{%
    \centering
    \texfigure[float=inline,max width=0.46\linewidth]{img_repaint/2014/shandong_liberal/q11_fig1.tex}%
}
\end{problem}

\begin{answer}
\(3\)
\end{answer}

\begin{solution}
初始时 \(x=1\)，\(n=0\). 判断条件为
\[
x^2-4x+3=(x-1)(x-3)\leqslant0
\]
即 \(1\leqslant x\leqslant3\). 当 \(x=1\) 时条件成立，执行一次循环后得到 \(x=2\)，\(n=1\)；当 \(x=2\) 时条件仍成立，执行一次循环后得到 \(x=3\)，\(n=2\)；当 \(x=3\) 时条件仍成立，执行一次循环后得到 \(x=4\)，\(n=3\). 此时 \(x^2-4x+3=3>0\)，循环结束，输出 \(n=3\).
\end{solution}

\begin{problem}
函数 \( y = \frac{\sqrt{3}}{2} \sin 2x + \cos^2 x \) 最小正周期为\fillinblank{}.
\end{problem}

\begin{answer}
\(\pi\)
\end{answer}

\begin{solution}
利用二倍角公式 \(\cos^2x=\frac{1+\cos2x}{2}\)，得
\[
y=\frac12+\frac{\sqrt3}{2}\sin2x+\frac12\cos2x
 =\frac12+\sin\left(2x+\frac\pi6\right)
\]
后式中正弦函数的自变量系数为 \(2\)，所以其最小正周期为 \(\frac{2\pi}{2}=\pi\). 因此原函数的最小正周期为 \(\pi\).
\end{solution}

\begin{problem}
一个六棱锥的体积为 \(2\sqrt{3}\)，其底面是边长为 2 的正六边形，侧棱长都相等，则该六棱锥的侧面积为\fillinblank{}.
\end{problem}

\begin{answer}
\(12\)
\end{answer}

\begin{solution}
设六棱锥的高为 \(h\). 正六边形底面的面积为
\[
S_\text{底}=6\times\frac{1}{2}\times2\times\sqrt3=6\sqrt3
\]
由棱锥体积公式，
\[
2\sqrt3=\frac13\times6\sqrt3\times h
\]
解得 \(h=1\). 由于六条侧棱相等，顶点在底面内的射影为正六边形的中心，设正六边形一边的中点为 \(E\)，则
\(OE=2\cos\frac\pi6=\sqrt3\)，\(PE=\sqrt{PO^2+OE^2}=\sqrt{1+3}=2\).
每个侧面三角形的面积为 \(\frac12\times2\times2=2\)，所以六个侧面的总面积为 \(6\times2=12\).
\end{solution}

\begin{problem}
圆心在直线 \(x - 2y = 0\) 上的圆 \(C\) 与 \(y\) 轴的正半轴相切，圆 \(C\) 截 \(x\) 轴所得弦的长为 \(2\sqrt{3}\)，则圆 \(C\) 的标准方程为\fillinblank{}.
\end{problem}

\begin{answer}
\((x - 2)^2 + (y - 1)^2 = 4\)
\end{answer}

\begin{solution}
设圆心为 \(O(h,k)\)，半径为 \(r\). 由圆心在直线 \(x-2y=0\) 上，得 \(h=2k\). 圆与 \(y\) 轴的正半轴相切，所以切点的纵坐标为正，且 \(k>0\)，半径为圆心到 \(y\) 轴的距离，即 \(r=h=2k\).

圆截 \(x\) 轴所得弦的长为 \(2\sqrt3\)，圆心到 \(x\) 轴的距离为 \(k\)，因此
\[
2\sqrt{r^2-k^2}=2\sqrt3
\]
代入 \(r=2k\)，得 \(4k^2-k^2=3\)，所以 \(k=1\)，从而 \(h=2\)，\(r=2\). 因此圆的标准方程为
\[
(x-2)^2+(y-1)^2=4
\]
\end{solution}

\begin{problem}
已知双曲线 \(\frac{x^2}{a^2} - \frac{y^2}{b^2} = 1\)（\(a > 0\)，\(b > 0\)）的焦距为 \(2c\)，右顶点为 \(A\)，抛物线 \(x^2 = 2py\)（\(p > 0\)）的焦点为 \(F\). 若双曲线截抛物线的准线所得线段长为 \(2c\)，且 \(FA = c\)，则双曲线的渐近线方程为\fillinblank{}.
\end{problem}

\begin{answer}
\(y = \pm x\)
\end{answer}

\begin{solution}
设双曲线的半焦距为 \(c\)，则
\[
c^2=a^2+b^2
\]
抛物线 \(x^2=2py\) 的准线为 \(y=-\frac p2\). 设双曲线与该准线的交点横坐标的绝对值为 \(x_0\)，则由题意 \(2x_0=2c\)，即 \(x_0=c\). 把 \(y=-\frac p2\) 代入双曲线方程，得
\[
\frac{x_0^2}{a^2}-\frac{p^2}{4b^2}=1
\]
利用 \(x_0=c\) 和 \(c^2=a^2+b^2\)，可得
\(\frac{a^2+b^2}{a^2}-\frac{p^2}{4b^2}=1\)，\(\frac{a^2p^2}{4b^2}=b^2\).
因为 \(a>0\)，\(b>0\)，\(p>0\)，所以 \(p=\frac{2b^2}{a}\). 抛物线焦点为 \(F\left(0,\frac p2\right)=\left(0,\frac{b^2}{a}\right)\)，双曲线右顶点为 \(A(a,0)\). 由 \(FA=c\)，得
\[
a^2+\frac{b^4}{a^2}=c^2=a^2+b^2
\]
于是 \(b^4=a^2b^2\)，结合 \(a>0\)，\(b>0\)，得 \(a=b\). 双曲线的渐近线为
\[
y=\pm\frac ba x=\pm x
\]
\end{solution}

\section{解答题}

\begin{problem}
海关对同时从 \(A\)，\(B\)，\(C\) 三个不同地区进口的某种商品进行抽样检测，从各地区进口此商品的数量（单位：\(\text{件}\)）如下表所示. 工作人员用分层抽样的方法从这些商品中共抽取 \(6\text{ 件}\)样品进行检测.

\begin{center}
\begin{tabular}{c|c|c|c}
\hline
地区 & \(A\) & \(B\) & \(C\) \\
\hline
数量 & 50 & 150 & 100 \\
\hline
\end{tabular}
\end{center}

\begin{enumerate}
\item 求这 \(6\text{ 件}\)样品中来自 \(A\)，\(B\)，\(C\) 各地区商品的数量；
\item 若在这 \(6\text{ 件}\)样品中随机抽取 \(2\text{ 件}\)送往甲机构进行进一步检测，求这 \(2\text{ 件}\)商品来自相同地区的概率.
\end{enumerate}
\end{problem}

\begin{solution}
\begin{enumerate}
\item 商品总数为 \(50+150+100=300\text{ 件}\)，抽取样品的比例为 \(6/300=1/50\). 因此，来自 \(A\)，\(B\)，\(C\) 三个地区的样品数分别为
\(50\times\frac1{50}=1\)，\(150\times\frac1{50}=3\)，\(100\times\frac1{50}=2\).
所以这 \(6\) 件样品中来自 \(A\)，\(B\)，\(C\) 三个地区的数量分别为 \(1\)，\(3\)，\(2\).

\item 从 \(6\) 件样品中任取 \(2\) 件，共有 \(\mathrm C_{6}^{2}\) 种取法. 若两件样品来自相同地区，则它们只能都来自 \(B\) 地区或都来自 \(C\) 地区，取法数为 \(\mathrm C_{3}^{2}+\mathrm C_{2}^{2}\). 因而所求概率为
\[
\frac{\mathrm C_{3}^{2}+\mathrm C_{2}^{2}}{\mathrm C_{6}^{2}}=\frac{3+1}{15}=\frac4{15}
\]
这里从 \(6\) 件样品中任取 \(2\) 件时每一种二件组合等可能，因此上述有利取法数与总取法数之比就是所求概率.
\end{enumerate}
\end{solution}

\begin{problem}
\(\triangle ABC\) 中，角 \(A\)，\(B\)，\(C\) 所对的边分别为 \(a\)，\(b\)，\(c\). 已知 \(a = 3\)，\(\cos A = \frac{\sqrt{6}}{3}\)，\(B = A + \frac{\pi}{2}\).

\begin{enumerate}
\item 求 \(b\) 的值；
\item 求 \(\triangle ABC\) 的面积.
\end{enumerate}
\end{problem}

\begin{solution}
由 \(A\) 为三角形内角，得 \(\sin A>0\)，所以
\[
\sin A=\sqrt{1-\cos^{2}A}=\sqrt{1-\frac{6}{9}}=\frac{\sqrt{3}}{3}
\]
又由 \(A+B+C=\pi\) 和 \(B=A+\frac\pi2\)，得 \(2A<\frac\pi2\)，从而 \(C=\frac\pi2-2A>0\)，这些角确实满足三角形内角的范围.

\begin{enumerate}
\item 由 \(B=A+\frac\pi2\)，得 \(\sin B=\cos A=\frac{\sqrt6}{3}\). 根据正弦定理，
\[
\frac{b}{\sin B}=\frac{a}{\sin A}
\]
从而
\[
b=a\frac{\sin B}{\sin A}=3\cdot\frac{\sqrt6/3}{\sqrt3/3}=3\sqrt2
\]

\item 三角形内角和给出 \(C=\pi-A-B=\frac\pi2-2A\)，于是
\[
\sin C=\cos 2A=\cos^{2}A-\sin^{2}A=\frac69-\frac39=\frac13
\]
因此三角形面积为
\[
S_{\triangle ABC}=\frac12ab\sin C=\frac12\cdot3\cdot3\sqrt2\cdot\frac13=\frac{3\sqrt2}{2}
\]
\end{enumerate}
\end{solution}

\begin{problem}
如图，四棱锥 \(P - ABCD\) 中，\(AP \perp\) 平面 \(PCD\)，\(AD \parallel BC\)，\(AB = BC = \frac{1}{2} AD\)，\(E\)，\(F\) 分别为线段 \(AD\)，\(PC\) 的中点.

\begin{enumerate}
\item 求证：\(AP \parallel\) 平面 \(BEF\)；
\item 求证：\(BE \perp\) 平面 \(PAC\).
\end{enumerate}

\bitmapfigure[max width=7.7cm]{img/2014/shandong_liberal/q18_fig1.png}
\end{problem}

\begin{solution}
设 \(\overrightarrow{BA}=\bs{a}\)，\(\overrightarrow{BC}=\bs{c}\)，\(\overrightarrow{BP}=\bs{p}\). 由四边形的顶点顺序、\(AD\parallel BC\) 以及 \(AD=2BC\)，有
\(\overrightarrow{AD}=2\overrightarrow{BC}\)，\(\overrightarrow{BD}=\bs{a}+2\bs{c}\).
由于 \(E\)，\(F\) 分别是 \(AD\)，\(PC\) 的中点，故
\(\overrightarrow{BE}=\bs{a}+\bs{c}\)，\(\overrightarrow{BF}=\frac{\bs{p}+\bs{c}}{2}\).

\begin{enumerate}
\item 由上式可得
\[
2\overrightarrow{BF}-\overrightarrow{BE}=\bs{p}+\bs{c}-\bs{a}-\bs{c}=\bs{p}-\bs{a}=\overrightarrow{AP}
\]
因此，直线 \(AP\) 的方向向量是平面 \(BEF\) 内两条相交直线 \(BE\)，\(BF\) 的方向向量的线性组合.

还需说明 \(AP\) 不在平面 \(BEF\) 内. 若 \(A\) 在平面 \(BEF\) 内，则存在实数 \(\alpha\)，\(\beta\)，使
\[
\bs{a}=\alpha(\bs{a}+\bs{c})+\frac{\beta}{2}(\bs{p}+\bs{c})
\]
由于 \(\bs{a}\)，\(\bs{c}\) 在底面内而 \(\bs{p}\) 不在底面内，比较垂直于底面的分量得 \(\beta=0\)，于是 \(\bs{a}=\alpha(\bs{a}+\bs{c})\)，这与四边形 \(ABCD\) 的底面不退化矛盾. 所以 \(AP\) 不在平面 \(BEF\) 内，结合其方向向量平行于该平面，得 \(AP\parallel\) 平面 \(BEF\).

\item 因为
\[
\overrightarrow{CD}=\overrightarrow{BD}-\overrightarrow{BC}=\bs{a}+\bs{c}=\overrightarrow{BE}
\]
所以 \(BE\parallel CD\). 又因为 \(CD\subset\) 平面 \(PCD\)，且 \(AP\perp\) 平面 \(PCD\)，故 \(BE\perp AP\).

另一方面，\(AB=BC\) 给出 \(\lvert\bs{a}\rvert=\lvert\bs{c}\rvert\). 而
\(\overrightarrow{AC}=\bs{c}-\bs{a}\)，\(\overrightarrow{BE}\cdot\overrightarrow{AC}=(\bs{a}+\bs{c})\cdot(\bs{c}-\bs{a})=\lvert\bs{c}\rvert^{2}-\lvert\bs{a}\rvert^{2}=0\)
所以 \(BE\perp AC\). 直线 \(AP\)，\(AC\) 是平面 \(PAC\) 内相交的两条直线，且 \(BE\) 同时垂直于它们，因此 \(BE\perp\) 平面 \(PAC\).
\end{enumerate}
\end{solution}

\begin{problem}
在等差数列 \(\{a_{n}\}\) 中，已知公差 \(d = 2\)，\(a_{2}\) 是 \(a_{1}\) 与 \(a_{4}\) 的等比中项.

\begin{enumerate}
\item 求数列 \(\{a_{n}\}\) 的通项公式；
\item 设 \(b_n = a_{\frac{n(n+1)}{2}}\)，记 \(T_n = -b_1 + b_2 - b_3 + b_4 - \cdots + (-1)^n b_n\)，求 \(T_n\).
\end{enumerate}
\end{problem}

\begin{solution}
\begin{enumerate}
\item 由公差 \(d=2\)，得
\(a_{2}=a_{1}+2\)，\(a_{4}=a_{1}+6\).
因为 \(a_{2}\) 是 \(a_{1}\) 与 \(a_{4}\) 的等比中项，所以
\[
(a_{1}+2)^{2}=a_{1}(a_{1}+6)
\]
展开并化简得 \(a_{1}^{2}+4a_{1}+4=a_{1}^{2}+6a_{1}\)，即 \(a_{1}=2\). 因此数列的通项公式为
\[
a_{n}=a_{1}+(n-1)d=2+2(n-1)=2n
\]

\item 由定义
\[
b_{n}=a_{\frac{n(n+1)}2}=2\cdot\frac{n(n+1)}2=n(n+1)
\]
当 \(n=2m\) 时，成对相加可得
\[
T_{2m}=\sum_{j=1}^{m}\bigl[-(2j-1)(2j)+(2j)(2j+1)\bigr]
=\sum_{j=1}^{m}4j=2m(m+1)
\]
当 \(n=2m+1\) 时，
\[
T_{2m+1}=T_{2m}-(2m+1)(2m+2)=-2(m+1)^{2}
\]
所以
\[
T_{n}=\begin{cases}
\dfrac{n(n+2)}2,&n=2m,\\[6pt]
-\dfrac{(n+1)^{2}}2,&n=2m+1,
\end{cases}
\]
其中偶数情形 \(m\geqslant1\)，奇数情形 \(m\geqslant0\).
\end{enumerate}
\end{solution}

\begin{problem}
设函数 \( f(x) = a \ln x + \frac{x - 1}{x + 1} \)，其中 \( a \) 为常数.

\begin{enumerate}
\item 若 \(a = 0\)，求曲线 \(y = f(x)\) 在点 \((1, f(1))\) 处的切线方程；
\item 讨论函数 \(f(x)\) 的单调性.
\end{enumerate}
\end{problem}

\begin{solution}
函数的定义域为 \(x>0\)，且
\[
f'(x)=\frac{a}{x}+\frac{2}{(x+1)^{2}}=\frac{a(x+1)^{2}+2x}{x(x+1)^{2}}
\]

\begin{enumerate}
\item 当 \(a=0\) 时，\(f(1)=0\)，\(f'(1)=\frac12\). 因而曲线在 \((1,f(1))=(1,0)\) 处的切线方程为
\[
y=\frac12(x-1)
\]

\item 由于 \(x(x+1)^{2}>0\)，\(f'(x)\) 的符号由
\[
g(x)=a(x+1)^{2}+2x=ax^{2}+2(a+1)x+a
\]
决定.

当 \(a\geqslant0\) 时，\(g(x)>0\)，所以 \(f(x)\) 在 \((0,+\infty)\) 上单调递增.

当 \(a<0\) 时，方程 \(g(x)=0\) 的判别式为 \(4(2a+1)\). 若 \(a<-\frac12\)，判别式小于零，且 \(g(0)=a<0\)，故 \(g(x)<0\)，函数在 \((0,+\infty)\) 上单调递减. 若 \(a=-\frac12\)，则
\[
g(x)=-\frac12(x-1)^{2}\leqslant0
\]
且 \(g(x)=0\) 仅在 \(x=1\) 处成立. 因此 \(f'(x)<0\) 在 \((0,1)\) 和 \((1,+\infty)\) 上均成立，函数在这两个区间上分别严格递减；又 \(f\) 在 \(x=1\) 处连续，所以对任意 \(0<x_1<1<x_2\)，有 \(f(x_1)>f(1)>f(x_2)\). 故函数在 \((0,+\infty)\) 上严格递减.

若 \(-\frac12<a<0\)，令
\(x_{1}=\frac{-(a+1)+\sqrt{2a+1}}{a}\)，\(x_{2}=\frac{-(a+1)-\sqrt{2a+1}}{a}\).
由于 \(a<0\) 且 \(0<\sqrt{2a+1}<a+1\)，有 \(0<x_{1}<x_{2}\). 此时 \(g(x)=0\) 的两个根为 \(x_{1}\)，\(x_{2}\)，又因二次项系数 \(a<0\)，故
\(f'(x)<0\)，\((0<x<x_{1}\text{ 或 }x>x_{2})\)，\(f'(x)>0\)，\((x_{1}<x<x_{2})\).
因此函数在 \((0,x_{1})\)，\((x_{2},+\infty)\) 上单调递减，在 \((x_{1},x_{2})\) 上单调递增.
\end{enumerate}
\end{solution}

\begin{problem}
在平面直角坐标系 \(xOy\) 中，椭圆 \(C: \frac{x^2}{a^2} + \frac{y^2}{b^2} = 1\)（\(a > b > 0\)）的离心率为 \(\frac{\sqrt{3}}{2}\)，直线 \(y = x\) 被椭圆 \(C\) 截得的线段长为 \(\frac{4\sqrt{10}}{5}\).

\begin{enumerate}
\item 求椭圆 \(C\) 的方程；
\item 过原点的直线与椭圆 \(C\) 交于 \(A\)，\(B\) 两点（\(A\)，\(B\) 不是椭圆 \(C\) 的顶点）. 点 \(D\) 在椭圆 \(C\) 上，且 \(AD \perp AB\)，直线 \(BD\) 与 \(x\) 轴，\(y\) 轴分别交于 \(M\)，\(N\) 两点.
\begin{enumerate}
\item 设直线 \(BD\)，\(AM\) 的斜率分别为 \(k_{1}\)，\(k_{2}\)，证明存在常数 \(\lambda\) 使得 \(k_{1} = \lambda k_{2}\)，并求出 \(\lambda\) 的值；
\item 求 \(\triangle OMN\) 面积的最大值.
\end{enumerate}
\end{enumerate}
\end{problem}

\begin{solution}
\begin{enumerate}
\item 设椭圆的半焦距为 \(c\). 由离心率条件
\(\frac{c}{a}=\frac{\sqrt3}{2}\)，\(c^{2}=a^{2}-b^{2}\)
得 \(b^{2}=a^{2}/4\). 直线 \(y=x\) 与椭圆的交点满足
\[
\frac{x^{2}}{a^{2}}+\frac{x^{2}}{a^{2}/4}=1
\]
即 \(x^{2}=a^{2}/5\). 两交点为 \((a/\sqrt5,a/\sqrt5)\)，\((-a/\sqrt5,-a/\sqrt5)\)，所以截得线段长为
\[
2\sqrt{\left(\frac{2a}{\sqrt5}\right)^{2}}=\frac{2\sqrt2a}{\sqrt5}
\]
由题设 \(\frac{2\sqrt2a}{\sqrt5}=\frac{4\sqrt{10}}5\)，解得 \(a=2\)，从而 \(b=1\). 因此椭圆方程为
\[
\frac{x^{2}}4+y^{2}=1
\]

\item 由于 \(A\)，\(B\) 不是椭圆的顶点，直线 \(AB\) 既不垂直于 \(x\) 轴，也不平行于 \(x\) 轴. 设其斜率为 \(t\)，则 \(t\neq0\). 令 \(A=(u,tu)\)，则由 \(A\) 在椭圆上知
\(\frac{u^{2}}4+t^{2}u^{2}=1\)，\(u^{2}=\frac4{1+4t^{2}}\).
因为直线 \(AB\) 过原点，椭圆关于原点对称，所以 \(B=(-u,-tu)\).

直线 \(AD\perp AB\)，而 \(AB\) 的方向向量为 \((1,t)\)，故可设
\[
D=A+s(t,-1)=(u+st,tu-s)
\]
将其代入椭圆方程，并利用 \(A\) 在椭圆上，得
\[
\frac{(u+st)^{2}}4+(tu-s)^{2}=1
\]
从而
\[
s\left(-\frac32ut+\frac{t^{2}+4}{4}s\right)=0
\]
根 \(s=0\) 对应点 \(A\)，另一个交点对应
\[
s=\frac{6ut}{t^{2}+4}
\]
因此
\[
D=\left(\frac{u(7t^{2}+4)}{t^{2}+4},\frac{ut(t^{2}-2)}{t^{2}+4}\right)
\]

\begin{enumerate}
\item 由上式和 \(B=(-u,-tu)\)，有
\[
\overrightarrow{BD}=\frac{2u(t^{2}+1)}{t^{2}+4}(4,t)
\]
所以
\[
k_{1}=\frac{t}{4}
\]
直线 \(BD\) 的方程为 \(y+tu=\frac{t}{4}(x+u)\). 令 \(y=0\)，得 \(M=(3u,0)\)；令 \(x=0\)，得 \(N=(0,-3tu/4)\). 因而
\[
k_{2}=k_{AM}=\frac{0-tu}{3u-u}=-\frac{t}{2}
\]
于是
\[
k_{1}=-\frac12k_{2}
\]
故存在常数 \(\lambda=-\frac12\)，使得 \(k_{1}=\lambda k_{2}\).

\item 由 \(M=(3u,0)\)，\(N=(0,-3tu/4)\)，得
\[
S_{\triangle OMN}=\frac12\lvert3u\rvert\cdot\left\lvert-\frac{3tu}{4}\right\rvert
=\frac98\lvert t\rvert u^{2}=\frac92\cdot\frac{\lvert t\rvert}{1+4t^{2}}
\]
令 \(r=\lvert t\rvert>0\). 由 \(A\)，\(B\) 不是椭圆顶点知 \(t\) 不能为 \(0\)，所以 \(r\) 的取值为全部正数. 因为
\[
1+4r^{2}-4r=(2r-1)^{2}\geqslant0
\]
所以 \(1+4r^{2}\geqslant4r\)，从而
\[
S_{\triangle OMN}\leqslant\frac92\cdot\frac14=\frac98
\]
当 \(r=\frac12\)，即 \(\lvert t\rvert=\frac12\) 时取等号，此时直线 \(AB\) 不是坐标轴，满足题设限制. 因此三角形 \(OMN\) 面积的最大值为 \(\frac98\).
\end{enumerate}
\end{enumerate}
\end{solution}
